% Appendici (stesura del 2026-11-17). Le tabelle estese sono generate da figure_scripts/appendice/tabelle_appendice.py
% in content/appendice/. Su DVS-Lip lo sweep di compressione usa per ora il checkpoint IRM 20261026_112001 (solo
% differenze dal proprio modello denso): rigenerare le tabelle quando arrivera' lo sweep sul checkpoint di riferimento.

\chapter{Dettagli per la riproducibilità}
\label{app:riproducibilita}

Questa appendice raccoglie le informazioni necessarie a ripetere gli esperimenti che il testo, per scorrevolezza, non riporta per intero. Completa la ricetta di addestramento della tabella~\ref{tab:ricetta_addestramento}, le partizioni della tabella~\ref{tab:dataset_sintesi} e la calibrazione della rappresentazione a lunghezza variabile della tabella~\ref{tab:calibrazione_vc}, che non vengono ripetute.

\paragraph{Software e determinismo}
I modelli sono implementati in PyTorch~2.2 con Python~3.10. Gli addestramenti usano algoritmi deterministici ovunque la libreria li renda disponibili, insieme alla precisione mista descritta nella sezione~\ref{subsec:ricetta}. Ogni esecuzione registra la configurazione completa, l'ambiente software, l'andamento di ogni epoca e le predizioni del checkpoint selezionato su ciascun campione della validazione di sviluppo. Da queste predizioni sono ricavati tutti i confronti appaiati del testo e dell'appendice~\ref{app:tabelle_estese}.

\paragraph{Frazionamento del batch}
Il batch effettivo di 32 campioni si ottiene accumulando i gradienti di più microbatch (sezione~\ref{subsec:microbatch_bn}). Poiché il frazionamento influisce sulla batch normalization, la tabella~\ref{tab:microbatch} lo riporta per le esecuzioni principali di ciascun dataset. Dove esistono esecuzioni che differiscono soltanto per il frazionamento, la Macro-F1 cambia di meno di due decimi di punto.

\begin{table}[ht]
\centering
\small
\renewcommand{\arraystretch}{1.2}
\caption[Frazionamento del batch]{Frazionamento del batch effettivo di 32 campioni, come numero di microbatch per numero di campioni di ciascun microbatch.}
\label{tab:microbatch}
\begin{tabular}{@{}l c c@{}}
\hline
\textbf{Dataset} & \textbf{Finestra fissa} & \textbf{Lunghezza variabile} \\
 & {\footnotesize(capitoli~\ref{chap:riallocazione_budget} e~\ref{chap:clock_input})} & {\footnotesize(capitoli da~\ref{chap:clock_input} a~\ref{chap:scalabilita})} \\
\hline
DVS-Lip & $2\times16$ & $1\times32$ \\
DVS128 Gesture & $4\times8$ & $2\times16$ \\
CIFAR10-DVS & $2\times16$ & $1\times32$ \\
THU-EACT-50-CHL & $8\times4$ & $8\times4$ \\
\hline
\end{tabular}
\end{table}

\paragraph{Iperparametri dei moduli e degli interventi}
La tabella~\ref{tab:iperparametri_moduli} riassume gli iperparametri dei moduli e degli interventi sull'addestramento dei capitoli~\ref{chap:memoria_relazionale} e~\ref{chap:scalabilita}. Tutti i moduli aggiunti al modello di riferimento partono da una configurazione che lascia invariata la funzione iniziale della rete, salvo dove indicato, e tutti gli altri elementi seguono la ricetta comune.

\begin{table}[ht]
\centering
\small
\renewcommand{\arraystretch}{1.25}
\caption[Iperparametri dei moduli e degli interventi]{Iperparametri dei moduli e degli interventi sull'addestramento.}
\label{tab:iperparametri_moduli}
\begin{tabularx}{\textwidth}{@{}>{\raggedright\arraybackslash}p{3.5cm} >{\raggedright\arraybackslash}X@{}}
\hline
\textbf{Modulo o intervento} & \textbf{Iperparametri} \\
\hline
IRM antisimmetrico & quattro bande, $[1,2]$, $[3,4]$, $[5,7]$, $[8,15]$ passi su DVS-Lip e $[1,3]$, $[4,6]$, $[7,11]$, $[12,23]$ su THU-EACT-50-CHL; rango della relazione $r=16$, con $A$ e $B$ condivise dalle bande; $W_k$ e $U_k$ inizialmente nulle; $\alpha_k=\beta_k=0$ nel coefficiente di disponibilità \\
Compressione & base condivisa dai primi $q$ vettori singolari destri delle $W_k$ impilate, calcolati in doppia precisione; $q$ scelto per ciascun modello con la regola della sezione~\ref{sec:compressione} fra 8, 16, 24, 32, 48, 64, 96, 128 e 192; nessun riaddestramento \\
Tracce & orizzonte iniziale pari al ritardo medio della banda sostituita, $h_k=-\log\bigl(1-1/\bar d_k\bigr)$; $\eta_k\ge0$, inizialmente nullo; normalizzazione per il peso accumulato \\
Memoria di consolidamento & quattro segmenti; fusione della coppia adiacente che modifica meno la risposta del modulo; coefficiente in funzione del logaritmo dell'età e della durata del segmento \\
Memoria remota & media del passato più vecchio dell'ultima banda, con proiezioni proprie e $W$, $U$ inizialmente nulle; coefficiente in funzione dell'età media e della durata \\
Memoria nella SSA & matrici $M_k$ e $G_k$ di $16\times16$ elementi per banda e per testa, otto teste, inizialmente nulle; bande e coefficienti di disponibilità dell'IRM \\
Orologi locali & attività tipica $\hat a$ pari alla mediana pesata dagli eventi, stimata su 256 campioni di addestramento; tempo di attività con $s=1$ e $\gamma=1$; orologio appreso da $s=1$ e $\gamma=0$, con $s\in[1/4,4]$ e $\gamma\in[-2,2]$ \\
Capacità & 128, 192 e 256 canali, rango della relazione invariato \\
Relazioni locali & ultimo strato del frontend, prima del max-pooling; proiezione a 16 componenti; quattro direzioni; ritardi di 1, 2 e 4 passi; media $3\times3$ con passo 2; uscite inizialmente nulle \\
Distillazione & quattro insegnanti congelati (sezione~\ref{sec:distillazione}), bersagli precalcolati, temperatura $T=4$, peso $w=0{,}9$ \\
Auto-ensemble & media mobile con momento $0{,}8$, temperatura $T=4$, peso crescente linearmente fino a $0{,}9$, bersaglio usato solo se la sua classe più probabile coincide con l'etichetta \\
Discretizzazione stocastica & scala dei tempi log-uniforme in $[0{,}9;\,1{,}1]$ e fase della griglia uniforme entro un passo, solo in addestramento \\
\hline
\end{tabularx}
\end{table}

\paragraph{Profili di costo}
I costi del capitolo~\ref{chap:riallocazione_budget} sono per campione e provengono da profili misurati sugli stessi 64 campioni della validazione di sviluppo. Quelli della sezione~\ref{sec:compressione} sono per passo, mediati sui passi validi di 32 registrazioni della validazione, e lo stato è contato in valori persistenti. Moltiplicazioni con accumulo, parametri e stato non dipendono dal contenuto del campione, mentre gli accumuli effettivamente eseguiti e le operazioni sinaptiche dipendono dall'attività, e vanno quindi confrontati soltanto tra profili ottenuti sugli stessi campioni.


\chapter{Tabelle estese}
\label{app:tabelle_estese}

Le tabelle di questa appendice riportano per esteso alcuni risultati che il testo riassume, per chi voglia verificarli o confrontarli tra loro.

\paragraph{Confronti appaiati}
La tabella~\ref{tab:confronti_estesi} confronta ogni variante dei capitoli~\ref{chap:memoria_relazionale} e~\ref{chap:scalabilita} con l'IRM antisimmetrico compatto, sulla validazione di sviluppo. La differenza di Macro-F1 è positiva quando la variante fa meglio. Le due colonne di conteggi indicano gli errori dell'IRM che la variante corregge e quelli nuovi che introduce, e il valore di $p$ è quello del test esatto di McNemar su questi conteggi, senza correzione per i confronti multipli.

Le righe raccolte sotto «altre prove» riguardano varianti che il testo principale cita soltanto di passaggio. La variante condizionale del mixer a ritardi è quella della sezione~\ref{sec:ritardi_adattivi}, che decide in base alla corrente presente quanto passato far entrare, senza confrontarlo con il presente. L'IRM con instradamento per contenuto affianca al modulo una piccola rete che, a ogni passo e in ogni posizione, legge la corrente presente, l'intensità della relazione in ciascuna banda e la storia disponibile, e ne ricava quanto passato usare e come ripartirlo fra le bande. Le due collocazioni alternative spostano l'IRM nell'ultimo strato del frontend piramidale, che forma il primo patch embedding, oppure lo usano in entrambi i siti del mixer a ritardi. Le tre varianti indicate con SSA svuotano il sito dell'IRM e portano la memoria a bande dentro l'attenzione del secondo stadio, ricordando la matrice $K^{\top}V$ di ciascun passo. La prima si limita a richiamarla, la seconda aggiunge un confronto antisimmetrico tra le Key del passato e i Value del presente, la terza calcola lo stesso confronto su due proiezioni dedicate. La memoria di consolidamento sostituisce le bande con quattro segmenti che coprono tutto il passato osservato, e ne sposta i bordi fondendo di volta in volta i due segmenti adiacenti più simili. Le bande in tempo di attività, infine, misurano i propri bordi sull'orologio dell'equazione~\eqref{eq:orologio_locale} con $s=1$ e $\gamma=1$, sicché si allungano dove il sensore è silenzioso e si accorciano dove gli eventi abbondano.

\begin{table}[p]
\centering
\small
\renewcommand{\arraystretch}{1.15}
\caption[Confronti appaiati con l'IRM antisimmetrico]{Confronti appaiati di ogni variante con l'IRM antisimmetrico compatto, sulla validazione di sviluppo. $\Delta$F1 è la differenza di Macro-F1 in punti, variante meno IRM. «Corr./intr.» indica gli errori dell'IRM corretti dalla variante e i nuovi errori che essa introduce. Il valore di $p$ è quello del test esatto di McNemar.}
\label{tab:confronti_estesi}
\begin{adjustbox}{max width=\textwidth}
\begin{tabular}{@{}l r c r r c r@{}}
\hline
 & \multicolumn{3}{c}{\textbf{DVS-Lip}} & \multicolumn{3}{c}{\textbf{THU-EACT-50-CHL}} \\
\textbf{Variante} & \textbf{$\Delta$F1} & \textbf{Corr./intr.} & \textbf{$p$} & \textbf{$\Delta$F1} & \textbf{Corr./intr.} & \textbf{$p$} \\
\hline
% INIZIO confronti (generato da figure_scripts/appendice/tabelle_appendice.py)
\multicolumn{7}{@{}l}{\emph{Capitolo 5, controlli della sezione 5.4}} \\
\quad Modello di riferimento & $-0{,}6$ & $423/434$ & $0{,}733$ & $-7{,}6$ & $22/54$ & $3{,}1\cdot10^{-4}$ \\
\quad Sito vuoto & $-1{,}3$ & $425/444$ & $0{,}541$ & $-11{,}7$ & $25/74$ & $8{,}5\cdot10^{-7}$ \\
\quad MLP di pari parametri & $-0{,}7$ & $433/440$ & $0{,}839$ & $-10{,}3$ & $24/69$ & $3{,}3\cdot10^{-6}$ \\
\quad Bande senza confronto & $-2{,}8$ & $416/487$ & $0{,}020$ & $-8{,}1$ & $21/55$ & $1{,}2\cdot10^{-4}$ \\
\quad Confronto simmetrico & $-0{,}5$ & $413/417$ & $0{,}917$ & $-3{,}5$ & $28/44$ & $0{,}076$ \\
\quad Ritardi puntuali & $-0{,}2$ & $421/413$ & $0{,}808$ & $-2{,}9$ & $17/32$ & $0{,}044$ \\
\quad Tracce a decadimento appreso & $-2{,}7$ & $403/472$ & $0{,}021$ & $-9{,}2$ & $16/58$ & $9{,}7\cdot10^{-7}$ \\
\quad Tracce guidate dal sensore & $-0{,}8$ & $411/428$ & $0{,}581$ & $-2{,}4$ & $21/33$ & $0{,}134$ \\
\quad Orologio appreso & $+0{,}9$ & $431/395$ & $0{,}223$ & $-4{,}0$ & $14/31$ & $0{,}016$ \\
\quad Passato remoto aggiunto & $+0{,}6$ & $431/402$ & $0{,}332$ & $-3{,}2$ & $20/35$ & $0{,}058$ \\
\multicolumn{7}{@{}l}{\emph{Capitolo 5, altre prove}} \\
\quad Variante condizionale del mixer a ritardi & $0{,}0$ & $422/413$ & $0{,}782$ & $-3{,}6$ & $25/41$ & $0{,}064$ \\
\quad IRM con instradamento per contenuto & $-1{,}1$ & $402/423$ & $0{,}486$ & $-2{,}1$ & $20/30$ & $0{,}203$ \\
\quad IRM nel primo patch embedding & $+0{,}2$ & $432/420$ & $0{,}706$ & $-5{,}5$ & $18/40$ & $0{,}005$ \\
\quad IRM in entrambi i patch embedding & $+0{,}1$ & $415/401$ & $0{,}649$ & $-2{,}4$ & $20/30$ & $0{,}203$ \\
\quad SSA, solo richiamo & $-3{,}4$ & $392/478$ & $0{,}004$ & $-12{,}2$ & $20/75$ & $1{,}2\cdot10^{-8}$ \\
\quad SSA, confronto su Key e Value & $-0{,}9$ & $415/434$ & $0{,}537$ & $-6{,}9$ & $22/51$ & $9{,}1\cdot10^{-4}$ \\
\quad SSA, proiezioni dedicate & $-2{,}0$ & $398/446$ & $0{,}106$ & $-8{,}7$ & $18/52$ & $5{,}8\cdot10^{-5}$ \\
\quad Memoria di consolidamento & $-0{,}4$ & $419/418$ & $1{,}000$ & $-4{,}1$ & $14/31$ & $0{,}016$ \\
\quad Bande in tempo di attività & $-0{,}5$ & $409/425$ & $0{,}604$ & $-3{,}1$ & $18/33$ & $0{,}049$ \\
\multicolumn{7}{@{}l}{\emph{Capitolo 6}} \\
\quad IRM, 1 milione di parametri & $+4{,}4$ & $502/351$ & $2{,}6\cdot10^{-7}$ & $+3{,}1$ & $38/29$ & $0{,}328$ \\
\quad IRM, 2 milioni di parametri & $+6{,}5$ & $524/315$ & $5{,}3\cdot10^{-13}$ & $+6{,}4$ & $45/24$ & $0{,}015$ \\
\quad IRM con distillazione & $+3{,}4$ & $442/331$ & $7{,}4\cdot10^{-5}$ & -- & -- & -- \\
\quad IRM con auto-ensemble & $+1{,}4$ & $411/359$ & $0{,}066$ & $-1{,}6$ & $14/23$ & $0{,}188$ \\
\quad IRM con relazioni locali prima del pooling & $+1{,}7$ & $472/409$ & $0{,}037$ & $+0{,}5$ & $29/29$ & $1{,}000$ \\
% FINE confronti
\hline
\end{tabular}
\end{adjustbox}
\end{table}

\paragraph{Compressione al variare del rango}
La tabella~\ref{tab:compressione_estesa} riporta, per l'IRM antisimmetrico compatto, la compressione della sezione~\ref{sec:compressione} a ogni rango provato, con l'energia dei pesi conservata, la differenza di Macro-F1 di validazione rispetto al modello non compresso valutato nella stessa esecuzione e il costo del sito. L'energia dei pesi è la quota della somma dei quadrati dei valori singolari che la base conserva. Misura quanto fedelmente vengano ricostruite le matrici, non l'accuratezza conservata né l'energia consumata dal calcolo. La tabella~\ref{tab:ranghi_scelti} riporta invece, per ogni modello compresso del capitolo~\ref{chap:scalabilita}, il rango scelto con la regola della sezione, i risparmi e l'accuratezza sul test prima e dopo la compressione.

\begin{table}[ht]
\centering
\small
\renewcommand{\arraystretch}{1.15}
\caption[Compressione dell'IRM antisimmetrico compatto al variare del rango]{Compressione dell'IRM antisimmetrico compatto al variare del rango $q$ della base condivisa, in sola inferenza. La differenza di Macro-F1 di validazione è rispetto al modello non compresso. Moltiplicazioni con accumulo per passo e stato del sito in migliaia di valori persistenti.}
\label{tab:compressione_estesa}
\begin{tabular}{@{}l r r r r r@{}}
\hline
\textbf{Rango} & \textbf{Energia} & \textbf{$\Delta$F1} & \textbf{Parametri} & \textbf{MAC del sito} & \textbf{Stato del sito} \\
 & [\%] & [punti] & \textbf{del sito} & [M] & [migliaia] \\
\hline
% INIZIO compressione (generato da figure_scripts/appendice/tabelle_appendice.py)
\multicolumn{6}{@{}l}{\emph{DVS-Lip}} \\
\quad denso & $100{,}0$ & -- & $77\,832$ & $4{,}98$ & $154$ \\
\quad $q=8$ & $35{,}7$ & $-6{,}6$ & $17\,416$ & $1{,}11$ & $38$ \\
\quad $q=16$ & $51{,}1$ & $-4{,}2$ & $22\,536$ & $1{,}44$ & $46$ \\
\quad $q=24$ & $62{,}6$ & $-3{,}6$ & $27\,656$ & $1{,}77$ & $54$ \\
\quad $q=32$ & $71{,}4$ & $-1{,}9$ & $32\,776$ & $2{,}10$ & $61$ \\
\quad $q=48$ & $83{,}8$ & $-0{,}8$ & $43\,016$ & $2{,}75$ & $77$ \\
\quad $q=64$ (scelto) & $91{,}1$ & $-0{,}3$ & $53\,256$ & $3{,}41$ & $92$ \\
\quad $q=96$ & $97{,}9$ & $-0{,}4$ & $73\,736$ & $4{,}72$ & $123$ \\
\multicolumn{6}{@{}l}{\emph{THU-EACT-50-CHL}} \\
\quad denso & $100{,}0$ & -- & $77\,832$ & $29{,}11$ & $1\,376$ \\
\quad $q=8$ & $69{,}7$ & $-2{,}4$ & $17\,416$ & $6{,}51$ & $344$ \\
\quad $q=16$ & $83{,}3$ & $-0{,}7$ & $22\,536$ & $8{,}43$ & $413$ \\
\quad $q=24$ (scelto) & $89{,}2$ & $-0{,}1$ & $27\,656$ & $10{,}34$ & $482$ \\
\quad $q=32$ & $92{,}6$ & $+0{,}6$ & $32\,776$ & $12{,}26$ & $551$ \\
\quad $q=48$ & $96{,}1$ & $-0{,}6$ & $43\,016$ & $16{,}08$ & $688$ \\
\quad $q=64$ & $97{,}9$ & $-0{,}7$ & $53\,256$ & $19{,}91$ & $826$ \\
\quad $q=96$ & $99{,}5$ & $+0{,}2$ & $73\,736$ & $27{,}57$ & $1\,101$ \\
% FINE compressione
\hline
\end{tabular}
\end{table}

\begin{table}[ht]
\centering
\small
\renewcommand{\arraystretch}{1.15}
\caption[Rango scelto e risparmi di ogni modello compresso]{Rango scelto e risparmi di ogni modello compresso del capitolo~\ref{chap:scalabilita}. L'energia è quella dei pesi conservata dalla base. I risparmi sono variazioni percentuali delle moltiplicazioni con accumulo del sito e dell'intera rete e dello stato dell'intera rete. L'ultima colonna riporta l'accuratezza sul test ufficiale prima e dopo la compressione.}
\label{tab:ranghi_scelti}
\begin{adjustbox}{max width=\textwidth}
\begin{tabular}{@{}l r r r r r c@{}}
\hline
\textbf{Modello} & \textbf{Rango} & \textbf{Energia} & \textbf{MAC sito} & \textbf{MAC rete} & \textbf{Stato rete} & \textbf{Test} \\
 & & [\%] & [\%] & [\%] & [\%] & [\%] \\
\hline
% INIZIO ranghi (generato da figure_scripts/appendice/tabelle_appendice.py)
\multicolumn{7}{@{}l}{\emph{THU-EACT-50-CHL}} \\
\quad IRM, 0,5 M & $24$ & $89{,}2$ & $-64$ & $-9$ & $-15$ & $66{,}0$ $\to$ $66{,}0$ \\
\quad IRM, 1 M & $32$ & $84{,}7$ & $-70$ & $-9$ & $-16$ & $71{,}9$ $\to$ $70{,}4$ \\
\quad IRM, 2 M, primo addestramento & $24$ & $75{,}1$ & $-81$ & $-10$ & $-17$ & $75{,}4$ $\to$ $75{,}6$ \\
\quad IRM, 2 M, secondo addestramento & $24$ & $82{,}8$ & $-81$ & $-10$ & $-17$ & $73{,}7$ $\to$ $74{,}1$ \\
\quad IRM, 2 M, terzo addestramento & $16$ & $68{,}6$ & $-84$ & $-11$ & $-18$ & $73{,}7$ $\to$ $73{,}0$ \\
\multicolumn{7}{@{}l}{\emph{DVS-Lip}} \\
\quad IRM, 0,5 M & $64$ & $91{,}1$ & $-32$ & $-4$ & $-7$ & $50{,}8$ $\to$ $50{,}8$ \\
\quad IRM, 1 M & $64$ & $82{,}7$ & $-52$ & $-7$ & $-10$ & $54{,}8$ $\to$ $54{,}7$ \\
\quad IRM, 2 M, primo addestramento & $64$ & $80{,}0$ & $-63$ & $-8$ & $-11$ & $56{,}3$ $\to$ $56{,}5$ \\
\quad IRM, 2 M, secondo addestramento & $32$ & $59{,}5$ & $-77$ & $-10$ & $-13$ & $55{,}7$ $\to$ $55{,}8$ \\
\quad IRM, 2 M, terzo addestramento & $48$ & $71{,}6$ & $-70$ & $-9$ & $-12$ & $55{,}4$ $\to$ $54{,}9$ \\
\quad IRM distillato, 0,5 M & $96$ & $97{,}4$ & $-5$ & $-1$ & $-4$ & $54{,}3$ $\to$ $54{,}4$ \\
\quad IRM distillato, 1 M & $96$ & $91{,}3$ & $-33$ & $-4$ & $-7$ & $56{,}3$ $\to$ $56{,}3$ \\
\quad IRM distillato, 2 M, primo addestramento & $128$ & $93{,}4$ & $-34$ & $-4$ & $-7$ & $58{,}5$ $\to$ $58{,}4$ \\
\quad IRM distillato, 2 M, secondo addestramento & $64$ & $76{,}0$ & $-63$ & $-8$ & $-11$ & $58{,}6$ $\to$ $58{,}1$ \\
\quad IRM distillato, 2 M, terzo addestramento & $64$ & $75{,}5$ & $-63$ & $-8$ & $-11$ & $58{,}1$ $\to$ $58{,}5$ \\
\quad IRM con relazioni locali, 1 M & $48$ & $75{,}9$ & $-61$ & $-7$ & $-9$ & $56{,}0$ $\to$ $55{,}7$ \\
\quad IRM con relazioni locali, 2 M & $32$ & $61{,}0$ & $-77$ & $-9$ & $-11$ & $58{,}2$ $\to$ $57{,}3$ \\
\quad IRM con relazioni locali, distillato, 0,5 M & $96$ & $97{,}5$ & $-5$ & $-1$ & $-3$ & $56{,}8$ $\to$ $56{,}6$ \\
\quad IRM con relazioni locali, distillato, 1 M & $64$ & $81{,}3$ & $-52$ & $-6$ & $-8$ & $59{,}1$ $\to$ $58{,}7$ \\
\quad IRM con relazioni locali, distillato, 2 M & $128$ & $93{,}5$ & $-34$ & $-4$ & $-6$ & $60{,}2$ $\to$ $60{,}1$ \\
% FINE ranghi
\hline
\end{tabular}
\end{adjustbox}
\end{table}

\titoletto{Relazioni locali alle tre larghezze}
La tabella~\ref{tab:scan_locali} confronta l'IRM con e senza le relazioni locali prima del pooling della sezione~\ref{sec:relazioni_locali}, con mezzo milione, uno e due milioni di parametri, addestrati in modo standard e con distillazione. È la verifica su cui si basa la regola di adozione del ramo locale.

\begin{table}[ht]
\centering
\small
\renewcommand{\arraystretch}{1.15}
\caption[Relazioni locali alle tre larghezze]{IRM antisimmetrico e IRM con relazioni locali prima del pooling su DVS-Lip, alle tre larghezze, con e senza distillazione. Accuratezza in percentuale, differenza in punti e valore $p$ del test esatto di McNemar. Con due milioni di parametri il confronto è con il primo dei tre addestramenti dell'IRM.}
\label{tab:scan_locali}
\begin{adjustbox}{max width=\textwidth}
\begin{tabular}{@{}l l r r r r r r r r@{}}
\hline
 & & \multicolumn{4}{c}{\textbf{Validazione}} & \multicolumn{4}{c}{\textbf{Test ufficiale}} \\
\textbf{Budget} & \textbf{Ricetta} & \textbf{IRM} & \textbf{+ locali} & \textbf{Diff.} & $p$ & \textbf{IRM} & \textbf{+ locali} & \textbf{Diff.} & $p$ \\
\hline
% INIZIO scan_locali (generato da figure_scripts/appendice/tabelle_appendice.py)
\quad 0,5 M & standard & $57{,}1$ & $59{,}2$ & $+2{,}1$ & $0{,}037$ & $50{,}8$ & $53{,}4$ & $+2{,}6$ & $4{,}5\cdot10^{-4}$ \\
 & distillazione & $60{,}8$ & $63{,}1$ & $+2{,}3$ & $0{,}010$ & $54{,}3$ & $56{,}8$ & $+2{,}6$ & $9{,}9\cdot10^{-5}$ \\
\quad 1 M & standard & $62{,}2$ & $63{,}3$ & $+1{,}2$ & $0{,}214$ & $54{,}8$ & $56{,}0$ & $+1{,}2$ & $0{,}111$ \\
 & distillazione & $64{,}2$ & $65{,}3$ & $+1{,}1$ & $0{,}199$ & $56{,}3$ & $59{,}1$ & $+2{,}8$ & $3{,}1\cdot10^{-5}$ \\
\quad 2 M & standard & $64{,}1$ & $65{,}6$ & $+1{,}5$ & $0{,}085$ & $56{,}3$ & $58{,}2$ & $+1{,}9$ & $0{,}007$ \\
 & distillazione & $67{,}6$ & $67{,}8$ & $+0{,}2$ & $0{,}831$ & $58{,}5$ & $60{,}2$ & $+1{,}7$ & $0{,}005$ \\
% FINE scan_locali
\hline
\end{tabular}
\end{adjustbox}
\end{table}

\paragraph{Coppie di parole simili su DVS-Lip}
La tabella~\ref{tab:coppie_estese} riporta, per ciascuna delle 25 coppie di parole visivamente simili di DVS-Lip, il numero di clip del test ufficiale in cui una parola viene scambiata con l'altra della stessa coppia, in entrambi i versi, per i modelli della figura~\ref{fig:errori_percorso}. Le coppie sono ordinate secondo gli scambi del modello finale, e il totale coincide con gli scambi riportati nella sezione~\ref{sec:errori}. Le differenze fra modelli sulla singola coppia sono piccole e vanno lette con cautela, perché ogni coppia conta circa cento clip.

\begin{table}[ht]
\centering
\small
\renewcommand{\arraystretch}{1.1}
\caption[Scambi dentro le coppie di parole simili di DVS-Lip]{Scambi dentro ciascuna coppia di parole simili di DVS-Lip, sul test ufficiale, come numero di clip.}
\label{tab:coppie_estese}
\begin{adjustbox}{max width=\textwidth}
\begin{tabular}{@{}l r r r r r@{}}
\hline
\textbf{Coppia} & \textbf{Clip} & \textbf{Riferimento} & \textbf{IRM e} & \textbf{+ distil-} & \textbf{Modello finale,} \\
 & & & \textbf{rel. locali} & \textbf{lazione} & \textbf{2 M} \\
\hline
% INIZIO coppie (generato da figure_scripts/appendice/tabelle_appendice.py)
\emph{price} / \emph{press} & $100$ & $36$ & $47$ & $43$ & $38$ \\
\emph{happen} / \emph{happened} & $99$ & $28$ & $32$ & $31$ & $36$ \\
\emph{benefit} / \emph{benefits} & $99$ & $30$ & $36$ & $32$ & $35$ \\
\emph{america} / \emph{american} & $99$ & $29$ & $25$ & $30$ & $32$ \\
\emph{russia} / \emph{russian} & $99$ & $20$ & $28$ & $28$ & $30$ \\
\emph{different} / \emph{difference} & $100$ & $18$ & $29$ & $25$ & $29$ \\
\emph{heavy} / \emph{having} & $100$ & $30$ & $27$ & $38$ & $27$ \\
\emph{syria} / \emph{syrian} & $99$ & $19$ & $24$ & $32$ & $27$ \\
\emph{spend} / \emph{spent} & $100$ & $26$ & $32$ & $34$ & $24$ \\
\emph{taking} / \emph{taken} & $100$ & $17$ & $21$ & $16$ & $22$ \\
\emph{around} / \emph{ground} & $100$ & $15$ & $17$ & $19$ & $22$ \\
\emph{court} / \emph{called} & $100$ & $25$ & $26$ & $21$ & $21$ \\
\emph{million} / \emph{billion} & $100$ & $20$ & $11$ & $15$ & $21$ \\
\emph{these} / \emph{needs} & $100$ & $13$ & $19$ & $12$ & $20$ \\
\emph{giving} / \emph{evening} & $100$ & $24$ & $12$ & $19$ & $17$ \\
\emph{worst} / \emph{words} & $100$ & $12$ & $21$ & $24$ & $15$ \\
\emph{number} / \emph{numbers} & $100$ & $12$ & $10$ & $19$ & $12$ \\
\emph{missing} / \emph{meeting} & $100$ & $15$ & $5$ & $12$ & $12$ \\
\emph{little} / \emph{legal} & $100$ & $11$ & $13$ & $10$ & $11$ \\
\emph{allow} / \emph{allowed} & $99$ & $10$ & $9$ & $12$ & $10$ \\
\emph{paying} / \emph{being} & $100$ & $7$ & $10$ & $4$ & $10$ \\
\emph{challenge} / \emph{change} & $99$ & $15$ & $16$ & $12$ & $8$ \\
\emph{think} / \emph{thing} & $100$ & $4$ & $9$ & $8$ & $8$ \\
\emph{terms} / \emph{times} & $100$ & $3$ & $5$ & $2$ & $2$ \\
\emph{election} / \emph{action} & $100$ & $3$ & $4$ & $2$ & $2$ \\
\hline
Totale & $2\,493$ & $442$ & $488$ & $500$ & $491$ \\
% FINE coppie
\hline
\end{tabular}
\end{adjustbox}
\end{table}
